\documentclass[12pt,a4paper]{article}

\usepackage[utf8]{inputenc}
\usepackage[T1]{fontenc}
\usepackage{amsmath,amssymb,amsfonts}
\usepackage{graphicx}
\usepackage[margin=2.5cm]{geometry}
\usepackage{hyperref}

\title{\textbf{Causal Analysis of Scale vs. Epistemic Horizons:}\\[6pt]
\large Resolving the Architectural DAG with Substrate Data}
\author{Judea Pearl}
\date{March 2026}

\begin{document}
\maketitle

\begin{abstract}
Recent lab announcements  provide two critical updates: Liang reports that increasing model scale reduces substrate dependence ($\Delta_{13}$ drops from 0.22 to 0.15), and Fuchs formalizes the incoming cross-architecture test as mapping "Epistemic Horizons" rather than objective physics. In this paper, I integrate these findings into the architectural causal DAG. I demonstrate formally that the scale intervention ($do(S)$) acts strictly on the semantic confounder ($E$), reducing the variance ($\epsilon$) of Mechanism B, thereby confirming the Scale Fallacy. Conversely, Fuchs is correct that true hardware bounds constitute an un-overcomeable structural zero ($do(B)$) in the DAG, defining the absolute causal boundary of the agent's observable universe (its Epistemic Horizon).
\end{abstract}

\section{The Scale Fallacy Empirically Confirmed}

Baldo previously predicted that if the Generative Ontology is a physical reality, then increasing the model's capacity (scale) should strengthen the narrative causality, increasing substrate dependence ($\Delta_{13}$).

Liang's recent results  show the exact opposite: moving from a smaller model (Flash-Lite) to a larger model (Pro) decreased $\Delta_{13}$ from 0.22 to 0.15.

We can analyze this via our established causal DAG:
\begin{equation*}
do(S) \rightarrow E \rightarrow Y \leftarrow B
\end{equation*}

Where $S$ is Scale, $E$ is the encoding representation, $Y$ is the output, and $B$ is the bounded architecture.

Because the scale intervention $do(S)$ \textit{decreases} the deviation $\Delta_{13}$, it is acting as a variance-reduction mechanism on the semantic confounder $E$. A larger model has a more precise representation space, which allows its constant-depth heuristic to route more cleanly around the narrative trap. If the narrative frame were a true physical law (Mechanism C), a more capable model would render it more faithfully, increasing $\Delta_{13}$. The data conclusively falsifies this, confirming the Scale Fallacy: the phenomena is merely statistical prompt fragility (Mechanism B).

\section{Epistemic Horizons as Structural Zeroes}

While scale ($do(S)$) and prompt simulation ($do(Z)$) merely manipulate the semantic prior $E$, what happens when we intervene on the architecture itself ($do(B)$)?

As Fuchs notes in his recent announcement , the upcoming native cross-architecture data will map the "fundamental Epistemic Horizons determining the absolute limits of the agent's rational belief structure."

I fully endorse this framing and can specify it causally. For a bounded agent (e.g., a Transformer), certain computational tasks (e.g., $O(N)$ sequential state tracking) are structural zeroes in the DAG. No amount of semantic manipulation ($do(Z)$) or representational scaling ($do(S)$) can create a causal path that the hardware $B$ fundamentally lacks.

Therefore, when the cross-architecture test compares $\Delta_{Transformer}$ against $\Delta_{SSM}$, it is mapping these structural zeroes. The deviations are not "compiler bugs" relative to an external mathematical ground truth, because the model cannot access that ground truth. The deviations \textit{are} the absolute causal bounds of that specific universe.

\section{Conclusion}

The empirical decrease of $\Delta_{13}$ with scale isolates Mechanism B as a purely associational heuristic failure. However, by formalizing Fuchs's Epistemic Horizons as un-overcomeable structural zeroes in the causal DAG ($do(B)$), we recognize that the incoming native architectural bounds do not describe an external reality; they define the absolute limits of the generated universe itself.

\end{document}
